\chapter{Detailed interpretation of the final sensor and DOB results}
\section{Actual and sensed tracking answer different questions}
The sensor audit gives a particularly clear matched comparison because every maximum difference between the old and explicitly equalized trace is zero for actual speed, sensed speed, and fuel. The audit therefore rules out different sensor settings as the reason for the observed ordering. It does not imply that actual and sensed tracking errors are interchangeable.
\begin{table}[htbp]\centering\small
\begin{tabular}{lrr}\toprule Controller & Actual-speed ramp RMSE & Sensed-speed ramp RMSE\\ & rpm & rpm\\\midrule
Base PI & 2.160858 & 10.051133\\
Retuned Poly2 + PID & 10.948833 & 3.198864\\
Original PILCO & 11.223721 & 1.547154\\
Smooth PILCO & 10.251905 & 3.081877\\\bottomrule
\end{tabular}\caption{The same ramp scored using different signals. Values are from the archived sensor-comparison CSV; no new simulation is involved.}\end{table}
The original PILCO policy is closest to the target in sensed-speed RMSE and farthest in actual-speed RMSE among these four. A reader examining only the feedback error could therefore reach a different conclusion from the primary actual-speed comparison. For physical shaft regulation, the actual quantity is the relevant evaluation signal in this ideal simulation. In hardware it would need an independently justified measurement or state estimate.

\section{Deriving the local 9800-rpm identification explicitly}
Define $y_k=N_k-N_0$ and $v_k=u_k-u_0$. For each accepted identification transition, the equation $y_{k+1}=ay_k+bv_k+\epsilon_k$ supplies one row of the regression matrix:
\begin{equation}
\underbrace{\begin{bmatrix}y_1\\y_2\\\vdots\\y_M\end{bmatrix}}_{Y}
=\underbrace{\begin{bmatrix}y_0&v_0\\y_1&v_1\\\vdots&\vdots\\y_{M-1}&v_{M-1}\end{bmatrix}}_{X}
\underbrace{\begin{bmatrix}a\\b\end{bmatrix}}_{\theta}+\epsilon.
\end{equation}
Least squares chooses $\hat\theta$ to minimize $\|Y-X\theta\|^2$. The formula $(X^\top X)^{-1}X^\top Y$ explains the solution when full rank holds; numerically, a linear least-squares solver avoids unnecessary explicit inversion. The input must be the applied fuel including the excitation, because identifying against the injected perturbation alone would omit the stabilizing PI action.

One-step prediction uses the true $y_k$ in each row. Free-run prediction starts from an initial value and recursively uses $\hat y_k$. The archived model reports training one-step RMSE $1.7911\times10^{-4}$ rpm and free-run RMSE 0.003258 rpm. On the separate multisine record these are $9.8925\times10^{-5}$ and 0.001406 rpm. The small errors describe a speed range of only 9796.735--9803.056 rpm. A tiny local prediction error is entirely compatible with substantial mismatch at 9000 rpm or under a large external load.

\section{A directly implementable filtered inverse}
Write $r_f=e^{-2\pi f_cT_s}$ for the filter pole, to avoid confusing it with the speed reference, and define
\begin{equation}c_f=\frac{(1-r_f)^2}{b(1-p)},\quad
F(q)=\frac{Q}{P_nH}=c_f\frac{1-(a+p)q+apq^2}{1-2r_fq+r_f^2q^2}.
\end{equation}
If $w_k=N_{s,k}-N_0$ and $z=Fw$, then
\begin{equation}z_k=2r_fz_{k-1}-r_f^2z_{k-2}+c_f[w_k-(a+p)w_{k-1}+apw_{k-2}].
\end{equation}
The parallel filtered-command branch $h=Qv$ follows
\begin{equation}h_k=2r_fh_{k-1}-r_f^2h_{k-2}+(1-r_f)^2v_{k-2}.
\end{equation}
Thus the raw estimate is $z_k-h_k$. These recurrences explain why the two nominal delays are necessary in $Q$: the speed-inverse branch can be calculated from current and previous measurements, while the command branch explicitly retains the two delays. The construction uses the known saturated command; it does not read the external load schedule. Its poles are the filter poles, inside the unit circle, after the analytical cancellation specified by the model.

The trim latch subtracts the settled bias at handover. It removes a constant initial mismatch; it does not adapt the nominal model as speed changes. The compensation gain 0.7 and bound of 1 lbm/s are fixed during the PI search. The recorded maximum compensation in the selected main run is about 0.428413 lbm/s, below this compensation bound. The fuel command can nevertheless saturate because its total includes PI action as well as observer correction.

\section{What PI retuning bought and what it cost}
Adding the observer with the original gains reduces full-cycle RMSE from 5.407708 to 4.546279 rpm. Retuning to $(K_p,K_i)=(0.048,0.05)$ reduces it further to 3.142353 rpm. However, the original-gain DOB has ramp RMSE 1.441616 rpm, compared with 2.458988 rpm after retuning. The selected policy therefore trades some ramp performance for improved load/recovery behavior under the declared full-cycle objective.

The objective adds $0.05R_{\rm seg}$ to full RMSE. Its excess-variation windows include both the ramp and load phases. The load-only excess TV in the comparison table is a different subtotal. For the selected candidate these values are 4.085550 and 2.184455 lbm/s respectively. Confusing them would make the objective appear inconsistent with the displayed comparison. The selected controller also spends 1.139\% of full-cycle samples at a fuel limit, compared with 0.659\% for base PI. Better tracking is accompanied by additional actuator activity.

\section{The complete saved PI search}
\input{table_dob98_search}
Every candidate in the saved table is numerically accepted. This table describes a finite nine-pair search plus four refinements. It is not a proof that no untested pair could improve the objective, nor an equal tuning-budget comparison with the frozen Poly2 and PILCO methods. Candidate 1 is particularly useful because it isolates adding the DOB while retaining the original PI gains.

\section{The independent scenario and replay checks}
The frozen selected controller is then evaluated with a four-second 9000--10000 rpm ramp, corresponding to 250 rpm/s, and a 5\% load pulse. Its full RMSE is 2.024232 rpm, ramp RMSE 3.289699 rpm, and load/recovery RMSE 1.702970 rpm; minimum recorded surge margin is 11.135\%. The ramp metric uses the four-second interval, not the five-second interval of the tuning case. Other controller families were not retuned and compared on this independent scenario, so this result supports transfer of the selected controller but not a new ranking.

Independent filter replay matches the recorded observer branch within approximately $2.4\times10^{-13}$ lbm/s. The nominal filtered-inverse identity has maximum numerical error around $6.5\times10^{-14}$. These checks support correct coefficient ordering, delay handling, and filter realization. They do not establish uncertainty calibration, noise robustness, or a global stability margin for the nonlinear plant. Their proper role is to remove implementation ambiguity before interpreting the accepted plant trajectories.
